% --------- Awards ------------------------------------------

\section{Awards}

\begin{itemize}[nosep, noitemsep]

    \item University of Central Florida Scroll \& Quill Society (October 2021).
    % , bestowed in recognition of sustained and outstanding achievements in research and/or creative activities at UCF
    
    \item Best Abstract Award, The 2nd Annual Women in IoT (WiT) (Virtual) Workshop: AI on the Edge Warren B. Nelms Institute for the Connected World, University of Florida for "A Privacy-Preserving Approach for Human Daily Activities Prediction", by S. Zehtabian, S. Khodadadeh, L. Bölöni and D. Turgut. 
    
    \item Finalist in the XPrize Pandemic Response Challenge with the team Pandemic Wave Predictor - PAWP (S. Zehtabian, S. Khodadadeh and L. Bölöni). Highest daily ranking: 5th out of 336 registered teams.
    
    \item University of Central Florida Research Incentive Award, 2019.
    
    \item Kurzweil Best AGI Idea Prize 2014 for the paper ``Autobiography based prediction in a situated AGI agent'', by L. Bölöni at the Seventh Conf. of Artificial General Intelligence (AGI-2014).
    
    \item Best Paper Award for the paper ``IVE: improving the value of information in energy-constrained intruder tracking sensor networks'', by D. Turgut and L. Bölöni at the IEEE Int. Conf. on Communications (ICC-2013).
    
    \item NASA Software Award, August 2006 for the NASA Engineering Shuttle Telemetry Agent - NESTA.
    
    \item AAAI Deployed Application Award for the paper ``NESTA: NASA Engineering Shuttle Telemetry Agent'' by G.S. Semmel, S.R. Davis, K.W. Leucht, D.A. Rowe, K.E. Smith, and L. Bölöni at the AAAI-2005 conference, July 2005.
    
    \item Member of the Upsilon Pi Epsilon Computer Sciences Honor Society, Beta Chapter of Indiana.
    
    \item Scholarship offered by Lucent Technologies for the university year 1998-99.
    
    \item Scholarship of the Hungarian Academy of Science the university year 1994-95, spent at the Analogical and Neural Computing Laboratory at the Computers and Automation Institute of Hungarian Academy of Science, working in the domain of cellular neural networks.
    
    \item Romanian Republican Scholarship in university year 1992-93 (awarded for one student in the department).
    
    \item First prize at the Romanian National Programming Contest of the students in 1990 (individual competition) and the third prize in  1991 and 1992 (with the team of the university).
    
\end{itemize}
